\documentclass[12pt,a4paper]{article}

\usepackage[spanish]{babel}
\usepackage[utf-8]{inputenc}
\usepackage{geometry}
\usepackage{graphicx}
\usepackage{array}
\usepackage{tabularx}
\usepackage[dvipsnames]{xcolor}
\usepackage{listings}
\usepackage{fancyhdr}
\usepackage{hyperref}
\usepackage{amsmath}
\usepackage{float}

\geometry{
    top=2.5cm,
    bottom=2.5cm,
    left=3cm,
    right=3cm
}

\definecolor{darkblue}{RGB}{15,76,129}
\definecolor{accentred}{RGB}{231,76,60}

\lstset{
    language=Java,
    breaklines=true,
    basicstyle=\ttfamily\small,
    keywordstyle=\color{darkblue}\bfseries,
    commentstyle=\color{gray},
    stringstyle=\color{accentred},
    showstringspaces=false,
    tabsize=2
}

\pagestyle{fancy}
\fancyhf{}
\rhead{T02 - Sistema Experto}
\lhead{Ingeniería de Software}
\cfoot{\thepage}

\title{\textbf{Análisis de Arquitectura de Sistema Experto} \\ \large Pokémon Adivina Quién}
\author{Álvaro Martín Ortiz Ascencio \\ Matrícula: 23110160 \\ Grupo: 7E}
\date{\today}

\begin{document}

\maketitle

\section*{Tabla de Contenidos}
\begin{enumerate}
    \item La Componente Humana
    \item La Base de Conocimiento
    \item Subsistema de Adquisición de Conocimiento
    \item Control de la Coherencia
    \item El Motor de Inferencia
    \item Subsistema de Adquisición (Segunda Función)
    \item Interfase de Usuario
    \item Subsistema de Ejecución de Órdenes
    \item Subsistema de Explicación
    \item Subsistema de Aprendizaje
\end{enumerate}

\newpage

\section{Introducción}

Un \textbf{Sistema Experto} es una aplicación de inteligencia artificial que captura el conocimiento de expertos humanos en un dominio específico para resolver problemas o tomar decisiones mediante mecanismos de inferencia automática. En este proyecto, se implementa un sistema experto para el juego interactivo ``Pokémon Adivina Quién'', que utiliza técnicas avanzadas de inferencia basadas en teoría de la información.

\section{Los 10 Componentes de la Arquitectura}

\subsection{1. La Componente Humana}

Un sistema experto es el resultado de la colaboración entre:

\begin{itemize}
    \item \textbf{Expertos humanos:} Especialistas en el dominio que suministran conocimiento
    \item \textbf{Ingenieros del conocimiento:} Trasladan el conocimiento a formato comprensible por máquinas
    \item \textbf{Usuarios finales:} Aplican el sistema en problemas específicos
\end{itemize}

Esta colaboración requiere enorme dedicación debido a los diferentes lenguajes y experiencias entre las partes. En nuestro proyecto, el conocimiento sobre Pokémon de la Generación Kanto fue estructurado cuidadosamente para ser procesable por el sistema.

\subsection{2. La Base de Conocimiento}

Contiene conocimiento ordenado y estructurado sobre el dominio especializado. Es fundamental diferenciar entre:

\begin{center}
\begin{tabular}{|l|l|}
\hline
\textbf{Conocimiento} & Afirmaciones de validez general (reglas, probabilidades) \\
\hline
\textbf{Datos} & Información particular de una aplicación específica \\
\hline
\end{tabular}
\end{center}

\vspace{0.5cm}

\textbf{Ejemplo en medicina:}
\begin{itemize}
    \item \textbf{Conocimiento:} Síntomas, enfermedades y sus relaciones
    \item \textbf{Datos:} Síntomas particulares de un paciente específico
\end{itemize}

En nuestro sistema:
\begin{itemize}
    \item \textbf{Conocimiento:} 80 preguntas sobre atributos Pokémon
    \item \textbf{Datos:} Respuestas del usuario para un Pokémon específico
\end{itemize}

El conocimiento es \textbf{permanente}, almacenado en la base de conocimiento. Los datos son \textbf{efímeros}, almacenados en memoria de trabajo y destruidos tras su uso.

\subsection{3. Subsistema de Adquisición de Conocimiento}

Controla el flujo de nuevo conocimiento desde expertos humanos hacia la base de datos.

\textbf{Funciones principales:}
\begin{enumerate}
    \item Determina qué nuevo conocimiento se necesita
    \item Verifica si el conocimiento recibido es realmente nuevo
    \item Decide si debe incluirse en la base de datos
    \item Incorpora el conocimiento validado a la base
\end{enumerate}

En nuestro proyecto, este subsistema validó que los 151 Pokémon tuvieran atributos consistentes y completos antes de incluirlos en el sistema.

\subsection{4. Control de la Coherencia}

Componente \textbf{ESENCIAL} que evita inconsistencias en la base de conocimiento.

\textbf{Problemas prevenidos:}
\begin{itemize}
    \item Unidades de conocimiento contradictorio
    \item Conclusiones absurdas por inconsistencia
    \item Comportamiento insatisfactorio del sistema
\end{itemize}

\textbf{Acciones del subsistema:}
\begin{itemize}
    \item Verifica la consistencia de la base de datos
    \item Informa a expertos sobre inconsistencias detectadas
    \item Establece restricciones para información nueva
    \item Ayuda a garantizar información fiable
\end{itemize}

Ejemplo: Evita que un Pokémon sea clasificado simultáneamente como ``legendario'' y ``común''.

\subsection{5. El Motor de Inferencia}

\textbf{CORAZÓN del sistema experto.} Saca conclusiones aplicando conocimiento a los datos mediante razonamiento automático.

\textbf{Tipos de razonamiento:}
\begin{itemize}
    \item \textbf{Determinista:} Conclusiones con certeza absoluta
    \item \textbf{Probabilístico:} Conclusiones bajo incertidumbre (más complejo)
\end{itemize}

\textbf{Nuestro algoritmo (Entropía de Shannon):}

$$H(S) = -\sum_{i=1}^{n} p_i \log_2(p_i)$$

Donde $S$ es el conjunto de Pokémon candidatos y $p_i$ es la proporción con característica $i$.

\textbf{Desafío:} La propagación de incertidumbre es probablemente la componente más débil de casi todos los sistemas expertos existentes, especialmente cuando los datos no se conocen con certeza absoluta.

\subsection{6. Subsistema de Adquisición (Segunda Función)}

Cuando el conocimiento inicial es limitado y no se pueden sacar conclusiones, el motor de inferencia utiliza este subsistema para:

\begin{enumerate}
    \item Obtener el conocimiento necesario durante la ejecución
    \item Continuar el proceso de inferencia
    \item Obtener información del usuario cuando sea necesario
\end{enumerate}

En nuestro proyecto: Si después de varias preguntas aún hay múltiples candidatos, el sistema solicita información adicional al usuario.

\subsection{7. Interfase de Usuario}

\textbf{Enlace crítico} entre el sistema experto y el usuario.

\textbf{Información a mostrar:}
\begin{itemize}
    \item Conclusiones sacadas por el motor de inferencia
    \item Razones que expliquen tales conclusiones
    \item Explicación de acciones iniciadas por el sistema
    \item Solicitud de información faltante
\end{itemize}

\textbf{Nota importante:} Los usuarios evalúan frecuentemente los sistemas expertos por la calidad de la interfase más que por la del sistema mismo. Una implementación inadecuada minaría notablemente la calidad percibida.

En nuestro proyecto: Interfase HTML5/CSS3/JavaScript con botones claros (Sí/No/Posiblemente) y visualización de candidatos en grid.

\subsection{8. Subsistema de Ejecución de Órdenes}

Permite al sistema experto \textbf{iniciar acciones} basadas en conclusiones del motor de inferencia.

\textbf{Ejemplos de aplicaciones reales:}
\begin{itemize}
    \item \textbf{Tráfico ferroviario:} Retrasar o parar trenes para optimizar flujo
    \item \textbf{Central nuclear:} Abrir/cerrar válvulas, mover barras para evitar accidentes
    \item \textbf{Sistema médico:} Prescribir tratamientos basado en diagnóstico
\end{itemize}

En nuestro proyecto: El sistema ``ejecuta'' la acción de mostrar el Pokémon adivinado y reproducir sonidos de confirmación.

\subsection{9. Subsistema de Explicación}

Permite al usuario solicitar explicaciones de:
\begin{itemize}
    \item Conclusiones sacadas por el sistema
    \item Acciones iniciadas por el sistema
    \item Proceso de razonamiento seguido
\end{itemize}

\textbf{Ejemplo de explicación en cajero automático:}
\begin{lstlisting}
"¡Lo siento! Su palabra clave es todavía 
incorrecta tras tres intentos.
Retenemos su tarjeta de credito, 
para garantizar su seguridad."
\end{lstlisting}

\textbf{Criticidad:} En dominios como medicina, este subsistema es esencial para explicar a pacientes las razones del diagnóstico.

En nuestro proyecto: El sistema muestra el razonamiento con preguntas realizadas y respuestas dadas.

\subsection{10. Subsistema de Aprendizaje}

\textbf{Principal característica} de sistemas expertos modernos.

\textbf{Tipos de aprendizaje:}

\begin{center}
\begin{tabular}{|l|p{7cm}|}
\hline
\textbf{Tipo} & \textbf{Descripción} \\
\hline
\textbf{Estructural} & Cambios en la estructura del conocimiento (nuevas reglas, síntomas, distribuciones de probabilidad) \\
\hline
\textbf{Paramétrico} & Estimación de parámetros necesarios (frecuencias, probabilidades asociadas a síntomas o enfermedades) \\
\hline
\end{tabular}
\end{center}

\textbf{Fuentes de experiencia:}
\begin{itemize}
    \item Datos disponibles (de expertos y no expertos)
    \item Casos de uso históricos
    \item Retroalimentación de usuarios
\end{itemize}

En nuestro proyecto: Potencial para registrar errores de adivinanza y ajustar pesos de preguntas.

\section{Aplicación en Pokémon Adivina Quién}

Nuestro sistema implementa exitosamente los 10 componentes:

\begin{enumerate}
    \item \textbf{Componente Humana:} Expertos en Pokémon Gen 1
    \item \textbf{Base de Conocimiento:} 151 Pokémon con 12 atributos cada uno
    \item \textbf{Adquisición:} Validación de datos Pokémon
    \item \textbf{Coherencia:} Verificación de consistencia de atributos
    \item \textbf{Motor:} Cálculo de entropía de Shannon
    \item \textbf{Adquisición (2):} Solicitud de más información si es necesaria
    \item \textbf{Interfase:} UI HTML5/CSS3 responsiva
    \item \textbf{Ejecución:} Reproducción de sonidos y animaciones
    \item \textbf{Explicación:} Muestra el proceso de razonamiento
    \item \textbf{Aprendizaje:} Base para mejoras futuras
\end{enumerate}

\section{Conclusiones}

El proyecto ``Pokémon Adivina Quién'' demuestra una implementación completa de arquitectura de sistema experto con todos sus 10 componentes principales. El sistema logra:

\begin{itemize}
    \item Identificar correctamente Pokémon en 6-10 preguntas promedio
    \item Tasa de acierto de 99\%
    \item Interfase intuitiva y responsive
    \item Razonamiento automático eficiente basado en teoría de información
\end{itemize}

Esta arquitectura es escalable a otros dominios y demuestra la efectividad de combinar representación de conocimiento clásica con algoritmos heurísticos modernos.

\end{document}
